En la era digital, las redes sociales se han consolidado como un espacio central para la comunicación, el comercio y el intercambio de información. Paralelamente, se han convertido en una fuente relevante de datos sobre experiencias cotidianas de los usuarios, incluyendo situaciones relacionadas con estafas, engaños y fraudes financieros. A través de publicaciones, comentarios y conversaciones, las personas comparten alertas, denuncias y relatos personales sobre intentos de fraude, lo que convierte a estas plataformas en un entorno privilegiado para el análisis de este tipo de fenómenos.

El fraude financiero digital no se limita únicamente a la ejecución directa de estafas a través de medios tecnológicos, sino que también se manifiesta en el discurso social que las rodea. Los usuarios describen intentos de \textit{phishing}, suplantaciones de identidad, estafas románticas, fraudes en compras \textit{online}, esquemas de inversión en criptomonedas o falsas comunicaciones institucionales, generando un volumen significativo de información textual que puede ser analizada de forma sistemática. Desde una perspectiva criminológica, el fraude se caracteriza por la intencionalidad, el engaño y el aprovechamiento de oportunidades, elementos que quedan reflejados en el lenguaje utilizado tanto por las víctimas como por los propios estafadores.

Los datos recientes ilustran la magnitud del fenómeno. En Estados Unidos, el informe \textit{Consumer Sentinel Network Data Book 2024} de la \textit{Federal Trade Commission} (FTC) registró pérdidas notificadas por fraude por valor de más de 12.500 millones de dólares durante 2024, lo que supone un incremento del 25\% respecto al año anterior. Las estafas de inversión encabezaron las pérdidas con 5.700 millones de dólares, seguidas por las estafas de suplantación con 2.950 millones. Es especialmente revelador que las estafas originadas en redes sociales concentren las mayores pérdidas por método de contacto, con 1.900 millones de dólares en 2024, lo que confirma el papel de estas plataformas como vector principal de la cibercriminalidad económica.

En España, el panorama es igualmente preocupante. El balance de ciberseguridad 2025 publicado por el Instituto Nacional de Ciberseguridad (INCIBE) refleja una crisis exponencial: 122.223 incidentes gestionados, un 26\% más que en 2024, de los cuales 45.445 corresponden a fraude \textit{online} (cuatro de cada diez incidentes). El \textit{phishing} lidera esta categoría con 25.133 casos, mientras que la Línea 017 de ayuda ciudadana atendió 142.767 consultas (un 44,9\% más que el año anterior), saturada por la demanda de asesoramiento preventivo y reactivo frente al fraude.

Estos datos revelan una realidad inquietante: la detección del fraude no puede entenderse únicamente como la identificación de transacciones ilícitas, sino también como la capacidad de reconocer señales tempranas en los patrones de comunicación. Las publicaciones y comentarios de los usuarios en redes sociales ofrecen la posibilidad de observar cómo emergen nuevas tipologías de fraude antes de que estas se consoliden o sean recogidas por los sistemas institucionales de alerta. Por ello, el análisis del lenguaje social se presenta como una herramienta complementaria crítica para la vigilancia y prevención del fraude financiero.

Entre las metodologías disponibles, el Procesamiento del Lenguaje Natural (\textit{Natural Language Processing}, NLP) permite analizar grandes volúmenes de texto de forma automatizada, identificando regularidades lingüísticas, temáticas y semánticas. Este enfoque resulta especialmente adecuado en entornos digitales, donde el fraude se comunica principalmente mediante mensajes escritos, enlaces, descripciones de experiencias o advertencias compartidas por los propios usuarios. Los avances recientes en modelos de lenguaje basados en \textit{transformers}, el desarrollo de técnicas de \textit{topic modeling} neuronal como BERTopic, y la aparición de clasificadores \textit{zero-shot} ofrecen hoy un instrumental maduro para abordar este tipo de análisis con datos reales de redes sociales.

Con el objetivo de identificar patrones discursivos asociados al fraude financiero con potencial predictivo, este trabajo se basa en el análisis de publicaciones extraídas de X (Twitter) en el contexto estadounidense durante el año 2025. A partir de estos datos, se busca construir un modelo que permita detectar tipologías emergentes de fraude y evaluar su posible transferencia al contexto español.

La elección de Estados Unidos como país fuente se fundamenta en tres razones principales. En primer lugar, la dimensión y madurez de su ecosistema digital convierten a este país en el principal generador de datos para entrenamiento de modelos de detección de fraude a escala global, con una enorme cantidad de usuarios activos en redes sociales y una abundancia de menciones relacionadas con estafas. En segundo lugar, la disponibilidad de registros oficiales consolidados (\textit{Consumer Sentinel Network} de la FTC, IC3 del FBI, CFPB) permite contrastar los patrones discursivos detectados con tipologías de fraude verificadas institucionalmente, aportando validez externa al modelo. En tercer lugar, Estados Unidos experimenta fraudes altamente industrializados (\textit{Fraud-as-a-Service}) apalancados en identidades sintéticas, \textit{deepfakes} y automatización a gran escala, lo que lo convierte en un laboratorio natural de tipologías emergentes que suelen aparecer con varios meses de adelanto respecto al contexto europeo.

España, por su parte, se configura como contexto destino por su creciente exposición al fraude digital y por su relevancia estratégica dentro del marco normativo europeo (RGPD, DORA, NIS2). La adopción masiva de banca digital, medios de pago electrónicos y redes sociales, combinada con un ecosistema de vigilancia institucional en expansión pero aún saturado, convierten al país en un escenario idóneo para evaluar la utilidad práctica de observatorios algorítmicos basados en NLP. La comparación bilateral Estados Unidos--España permite así articular un puente metodológico entre un contexto anglosajón maduro en la generación de patrones de fraude y un contexto hispanohablante que requiere herramientas adaptadas para la interceptación del fraude en fase discursiva.

\subsection{Objetivos}

El objetivo principal del presente Trabajo Fin de Máster es desarrollar un modelo basado en Procesamiento del Lenguaje Natural (NLP) capaz de identificar y analizar menciones relacionadas con fraude financiero en redes sociales, con el fin de detectar patrones recurrentes y tipologías emergentes que puedan anticipar riesgos en el contexto español a partir de la evidencia observada en Estados Unidos.

Para alcanzar este objetivo general, se plantean los siguientes objetivos específicos:

\begin{itemize}
    \item Recopilar datos textuales procedentes de redes sociales mediante el uso de la API de X (Twitter), correspondientes a Estados Unidos durante el año 2025, utilizando consultas específicas y filtros geográficos que garanticen la relevancia del corpus.
    \item Preprocesar y limpiar los datos, eliminando ruido, duplicados, \textit{tweets} irrelevantes y contenido no pertinente, con el fin de obtener un corpus textual adecuado para el análisis lingüístico y semántico.
    \item Explorar la estructura temática del corpus mediante \textit{topic modeling} no supervisado con BERTopic, identificando los grandes bloques discursivos presentes en las menciones de fraude sin imponer categorías a priori.
    \item Clasificar automáticamente las menciones de fraude financiero en categorías definidas mediante técnicas de clasificación \textit{zero-shot} con modelos preentrenados, asignando a cada \textit{tweet} una tipología concreta (estafa de inversión, estafa romántica, \textit{phishing}, suplantación gubernamental, fraude en \textit{apps} de pago, etc.) o marcándolo como irrelevante.
    \item Analizar patrones lingüísticos y semánticos presentes en las conversaciones sobre fraude, con el objetivo de detectar estructuras discursivas recurrentes (urgencia, autoridad simulada, promesas de rentabilidad) asociadas a cada tipología.
    \item Traducir al español los contenidos clasificados mediante modelos de traducción neuronal (Helsinki-NLP Opus-MT), preservando el valor informativo original para su análisis desde el contexto español.
    \item Evaluar la transferibilidad del modelo al caso español, comprobando si los patrones detectados en Estados Unidos pueden utilizarse como indicadores tempranos de nuevas formas de fraude en España y qué ajustes requiere su traslación al ecosistema hispanohablante.
    \item Explorar el potencial de las redes sociales como herramienta de vigilancia temprana, aportando una metodología alternativa y complementaria a los sistemas institucionales de alerta (INCIBE, Policía Nacional, CNMV) basada en el análisis del lenguaje social.
\end{itemize}

\subsection{Estado del Arte}

\subsubsection{Fenomenología del Fraude y el Paradigma del Análisis Social}

La convergencia entre la digitalización financiera y la adopción masiva de redes sociales ha reconfigurado el panorama de la cibercriminalidad. Debido a la robustificación de la seguridad bancaria tradicional, los atacantes han desplazado su foco hacia el factor humano mediante técnicas de ingeniería social. En este contexto, plataformas como X (Twitter) o Reddit asumen una dualidad: son el vector principal para la ejecución de estafas y, simultáneamente, actúan como un repositorio público donde las víctimas denuncian y alertan sobre fraudes incipientes.

La detección transaccional tradicional adolece de una naturaleza reactiva: para cuando el sistema bancario detecta la anomalía, el daño patrimonial suele ser irreversible. Frente a esto, el Procesamiento de Lenguaje Natural (NLP) aplicado al discurso social permite una vigilancia proactiva. El fraude es un proceso persuasivo que deja una «huella semántica» observable mucho antes de materializarse económicamente. Extraer este conocimiento permite construir sistemas de alerta temprana basados en las narrativas de los usuarios, operando en el espacio temporal previo a la materialización del daño económico.

\subsubsection{Fundamentos Lingüísticos del Engaño}

El análisis de NLP requiere comprender cómo se manifiestan la coerción y el riesgo en el lenguaje. Los algoritmos se entrenan para ponderar características lingüísticas que correlacionan con la intención maliciosa. Entre las más relevantes se encuentran:

\begin{itemize}
    \item \textbf{Condicionalidad temporal (FOMO):} uso de límites de tiempo estrictos y adverbios de urgencia («última oportunidad», «solo hoy», «antes de medianoche») para forzar atajos heurísticos en la víctima y neutralizar el pensamiento deliberativo.
    \item \textbf{Prueba social sintética:} inclusión de jerga financiera descontextualizada, testimonios falsos o asociación ficticia con entidades reguladoras para manufacturar autoridad y credibilidad.
    \item \textbf{Anomalías morfosintácticas:} inconsistencias de registro, errores de concordancia o giros atípicos derivados de traducciones automáticas en campañas transnacionales, que actúan como marcadores estadísticos de origen fraudulento.
\end{itemize}

Paralelamente, el análisis de las narrativas de las víctimas actúa como un sensor de alta fidelidad. Los modelos de NLP que detectan emociones como ansiedad, miedo y culpa vinculadas a menciones de pasarelas de pago permiten generar alertas comunitarias (\textit{crowdsourced alerts}) mucho más ágiles que las listas negras corporativas.

\subsubsection{Evolución Arquitectónica: De Modelos Estadísticos a LLMs}

La literatura científica muestra una rápida evolución en las arquitecturas para procesar textos fraudulentos. En una primera etapa, los enfoques estadísticos e híbridos clásicos dominaron el campo. Metodologías como \textit{Bag of Words} y TF-IDF lideraron la detección inicial: el trabajo pionero de \cite{Dong16} integró un Modelo Estadístico de Lenguaje (SLM) con Análisis Semántico Latente (LSA) para capturar el uso de lenguaje engañoso en reportes corporativos y redes sociales, logrando precisiones superiores al 75\% usando SVM. Del mismo modo, estudios como el de \cite{Gangwar22} demostraron que combinar la semántica del texto con metadatos del usuario elevaba la precisión en X (Twitter) por encima del 90\%.

Para superar la «ceguera contextual» de los modelos estadísticos, investigadores como \cite{Li23} aplicaron \textit{Deep Learning} mediante \textit{embeddings} densos, aunque evidenciaron problemas de sobreajuste (\textit{overfitting}) debido al desbalance de clases, dado que el fraude constituye un evento estadísticamente raro. Para solucionar esto, \cite{Shukla25} propusieron una arquitectura híbrida de vanguardia: una Red Neuronal Convolucional Temporal (TCN) para asimilar secuencias de comportamiento a largo plazo, combinada con una Red Generativa Antagónica (GAN). La GAN sintetiza datos fraudulentos hiperrealistas para equilibrar el conjunto de entrenamiento, logrando un ROC-AUC de 0,96 en bases de datos sociales complejas.

El estado del arte actual pivota hacia los Grandes Modelos de Lenguaje (LLMs). \cite{Deng23} abordaron la complejidad y el ruido de plataformas como Reddit utilizando aprendizaje semi-supervisado con el modelo PaLM. Ante la falta de datos etiquetados, utilizaron razonamiento en Cadena de Pensamiento (CoT), forzando al LLM a resumir argumentos financieros antes de predecir. Emplearon el LLM masivo como «profesor» para generar puntuaciones suaves (\textit{soft scores}) y destilaron ese conocimiento en un modelo «estudiante» pequeño y eficiente para producción, igualando el rendimiento de los humanos expertos. Complementariamente, iniciativas como BeyondPhish (2023) demostraron que analizar semánticamente foros como r/Scams detecta el 98,3\% de fraudes de comercio electrónico, superando ampliamente a los motores corporativos tradicionales.

Dentro de este ecosistema técnico, BERTopic \citep{Grootendorst22} se ha consolidado como una de las técnicas de \textit{topic modeling} más utilizadas en investigación aplicada a redes sociales. Combina \textit{embeddings} densos generados por modelos tipo Sentence-BERT, reducción de dimensionalidad mediante UMAP, \textit{clustering} jerárquico con HDBSCAN y una variante de TF-IDF basada en clases (c-TF-IDF) para extraer términos representativos. Frente a métodos clásicos como LDA, BERTopic ofrece mayor coherencia semántica, mejor interpretabilidad y adaptación natural a textos cortos y ruidosos como los de X (Twitter). Su variante \textit{zero-shot} permite además definir tópicos predefinidos y asignar documentos a ellos por similitud coseno, combinando exploración no supervisada y clasificación dirigida en un único flujo.

\subsubsection{Transferibilidad Translingüística (Cross-Lingual Transfer)}

Para aplicar los patrones anglosajones detectados en Estados Unidos al contexto hispanohablante español, se requiere recurrir al aprendizaje de transferencia translingüística. Los modelos multilingües (mPLMs) mapean conceptos de coerción, urgencia o autoridad en un espacio geométrico unificado, independientemente del idioma, lo que permite reutilizar modelos entrenados en corpus anglosajones masivos para aplicar sobre datos en castellano.

El estudio «Dólares or Dollars?» \citep{Zhang24} introdujo FinMA-ES, el primer LLM optimizado de forma bilingüe para finanzas. Su investigación demostró un «efecto de transferencia lingüística cruzada beneficioso»: entrenar el modelo simultáneamente con datos en inglés y español no solo le permitió superar a GPT-4 en tareas en español, sino que también mejoró su precisión en inglés. Esto valida la viabilidad de transferir patrones de detección de estafas sin necesidad de reconstruir corpus masivos en castellano desde cero, un aspecto metodológicamente crítico para el presente trabajo, que parte de datos recolectados en inglés para su aplicación al contexto español.

\subsubsection{Análisis Comparativo por Jurisdicciones}

Aunque la táctica psicológica del engaño es esencialmente universal, las vías de explotación presentan diferencias estructurales entre ambos contextos que conviene caracterizar.

Estados Unidos actúa como principal generador de datos para entrenamiento a escala global. Experimenta fraudes altamente industrializados (\textit{Fraud-as-a-Service}) apalancados en identidades sintéticas y \textit{deepfakes}. Según el \textit{Consumer Sentinel Network Data Book 2024} de la FTC, las pérdidas declaradas superaron los 12.500 millones de dólares en 2024, con 2.600.000 reportes de fraude. Las estafas de inversión encabezan las pérdidas (5.700 millones, +24\% interanual), seguidas por las estafas de suplantación (2.950 millones). Especialmente relevantes para el presente trabajo son dos datos: primero, las estafas originadas en redes sociales concentran las mayores pérdidas por método de contacto (1.900 millones de dólares); segundo, las estafas de suplantación gubernamental se dispararon un 600\% hasta los 789 millones de dólares. Esta última tendencia es particularmente significativa como tipología emergente susceptible de transferirse al contexto europeo.

España, por su parte, presenta una situación caracterizada por un crecimiento acelerado del fraude digital en paralelo a la saturación de los mecanismos institucionales de respuesta. El balance de ciberseguridad 2025 del INCIBE refleja 122.223 incidentes gestionados (+26\% interanual), de los cuales 45.445 corresponden a fraude \textit{online} (+19\%). El \textit{phishing} lidera la categoría con 25.133 casos, y se cerraron 4.600 dominios .es potencialmente fraudulentos en colaboración con Red.es. La Línea 017 atendió 142.767 consultas (+44,9\%), saturada por la demanda ciudadana: un 28\% de los usuarios recibió intentos de \textit{phishing}, \textit{vishing} o \textit{smishing}, un 16\% solicitó ayuda tras compras fraudulentas \textit{online}, y un 14\% por suplantación de identidad digital. Ante este volumen masivo y una atención ciudadana desbordada, la implantación de observatorios algorítmicos basados en NLP resulta crítica para interceptar el fraude en fase discursiva y complementar los sistemas reactivos tradicionales.

\subsubsection{Limitaciones y Conclusión}

La aplicación de NLP en producción enfrenta desafíos éticos y técnicos de relevancia creciente. Destacan la opacidad algorítmica (caja negra), que choca con la regulación europea (RGPD, Reglamento de IA) sobre explicabilidad de la inteligencia artificial (XAI); la deriva semántica (\textit{concept drift}), ya que los estafadores mutan su vocabulario rápidamente para evadir la detección; y el riesgo de envenenamiento de datos, donde cibercriminales podrían inyectar textos manipulados en redes sociales para alterar los análisis de sentimiento financiero y la confianza del mercado. A ello se suma la tensión entre el uso de datos públicos de redes sociales para investigación y las restricciones cada vez más estrictas de las APIs comerciales, que condicionan la reproducibilidad de los estudios.

En conclusión, la metodología reactiva basada puramente en transacciones resulta insuficiente ante un fenómeno que opera a velocidad algorítmica. Las redes sociales se han consolidado como sensores tempranos formidables. El progreso desde técnicas de \textit{Bag of Words} hasta LLMs con Cadena de Pensamiento y destilación de conocimiento, junto con la demostrada viabilidad de la transferencia translingüística, confirma que España puede y debe capitalizar los patrones globales de fraude extraídos del ámbito anglosajón para la protección activa de su ecosistema financiero. El presente Trabajo Fin de Máster se inscribe en esta línea, aportando una aplicación empírica concreta del paradigma.
